\section{Dari Mnemonic ke Kode Mesin}
\label{sec:mnemonic-ke-kode-mesin}

Bab-bab sebelumnya telah memperkenalkan register dan segmentasi
sebagai wadah data. Bab ini melangkah satu tingkat lebih dalam: ke
bentuk fisik instruksi yang benar-benar dibaca prosesor. Pemahaman
atas format instruksi adalah kunci untuk membaca hasil
\eng{disassembly}, menghitung ukuran program, dan menelusuri kesalahan
pada level mesin.

\subsection{Assembly adalah Simbol dari Kode Mesin}

Perangkat keras hanya mengenal deretan bit. Instruksi untuk
memindahkan isi register \reg{BX} ke register \reg{AX} tidak pernah
berbentuk huruf \code{MOV}; yang tersimpan di memori adalah dua byte
\hex{89 D8}. Deretan byte inilah yang disebut \term{kode mesin}
(\eng{machine code}).

Manusia menulis \lstasm|MOV AX, BX| karena deretan byte jauh lebih
sulit diingat dan rawan salah tulis. Kata pendek yang mewakili operasi
tersebut disebut \term{mnemonic}. Assembler menerjemahkan setiap
mnemonic beserta operandnya menjadi deretan byte kode mesin yang
ekuivalen.

Pertanyaan pentingnya: mengapa \asm{MOV} tidak selalu menghasilkan
opcode yang sama? Jawabannya terletak pada operand. \asm{MOV} dapat
memindahkan data antardua register, dari nilai segera ke register,
dari memori ke register, atau sebaliknya. Setiap bentuk operand
memerlukan pengodean yang berbeda.

\subsection{Mengapa Satu Mnemonic Memiliki Beberapa Encoding}

Ada empat sebab utama yang membuat satu mnemonic memiliki banyak
bentuk \eng{encoding}.

\begin{enumerate}
    \item \textbf{Jenis operand.} Register, sel memori, dan nilai
          segera adalah tiga jenis operand yang berbeda. Instruksi
          \lstasm|MOV AX, BX| (register ke register) memakai opcode
          \hex{89}, sedangkan \lstasm|MOV AX, 1234h| (nilai segera ke
          register) memakai opcode \hex{B8}.

    \item \textbf{Ukuran operand.} Prosesor 8086 bekerja dengan data
          8 bit dan 16 bit. Instruksi \lstasm|MOV AL, 5| dan
          \lstasm|MOV AX, 5| beroperasi pada ukuran berbeda dan
          memakai opcode yang berbeda pula, yaitu \hex{B0} dan
          \hex{B8}.

    \item \textbf{Arah perpindahan.} Opcode \hex{8B} berarti
          ``register $\leftarrow$ memori'', sedangkan \hex{89}
          berarti ``memori $\leftarrow$ register''. Assembler memilih
          salah satunya sesuai urutan operand yang ditulis programmer.

    \item \textbf{Mode pengalamatan.} Operand memori dapat berupa
          alamat langsung, penjumlahan dua register, atau penjumlahan
          dengan sebuah konstanta. Perbedaan ini diatur melalui byte
          ModR/M dan \eng{displacement}, bukan dengan mengganti opcode
          (Bab~\ref{chap:mode-pengalamatan}).
\end{enumerate}

Akibatnya \asm{MOV} memiliki puluhan bentuk mesin. Programmer tidak
perlu menghafalnya karena assembler yang memilih bentuk yang tepat,
tetapi pembaca hasil \eng{disassembly} harus mengenali bentuk yang
sedang dihadapi.

\subsection{Format Umum Instruksi 8086}

Meskipun jumlah bentuknya banyak, seluruh instruksi 8086 mengikuti
satu pola susunan field yang sama. Panjang totalnya antara 1 sampai
6 byte dan tersusun dari lima bagian berikut, dengan urutan yang
tetap.

\begin{enumerate}
    \item \textbf{Byte awalan} (\eng{prefix}), 0--1 byte. Bersifat
          pilihan. Contohnya \code{LOCK} (\hex{F0}), pengganti segmen
          \code{ES} (\hex{26}), dan \code{REP} (\hex{F3}).

    \item \textbf{Opcode}, 1--2 byte. Menyatakan operasi yang harus
          dilakukan, misalnya ``pindahkan'', ``tambahkan'', atau
          ``lompat''.

    \item \textbf{Byte ModR/M}, 0--1 byte. Menyatakan letak operand
          dan cara pengalamatannya. Byte ini hadir bila instruksi
          melibatkan register atau memori.

    \item \textbf{Displacement}, 0--2 byte. Nilai pergeseran alamat
          pada operand memori berbentuk \lstasm|[SI+4]|.

    \item \textbf{Immediate}, 0--2 byte. Nilai konstanta yang menjadi
          operand langsung, misalnya \hex{5} pada \lstasm|MOV AL, 5|.
\end{enumerate}

Field yang tidak diperlukan tidak ditulis sama sekali, sehingga
panjang instruksi bervariasi. Susunan kelima field tersebut
ditunjukkan pada Gambar~\ref{fig:format-instruksi}.

\begin{figure}[H]
    \centering
    \begin{tikzpicture}[
        font=\small,
        bidang/.style={
            draw=obeBlue,
            line width=0.8pt,
            minimum height=1.0cm,
            align=center,
            inner sep=3pt
        },
        kosong/.style={
            draw=obeGray,
            line width=0.8pt,
            dashed,
            minimum height=1.0cm,
            align=center,
            inner sep=3pt
        },
        ket/.style={font=\scriptsize, text=obeGray, align=center}
    ]
        % --- baris pertama: kerangka umum ---
        \node[bidang, fill=obeAmberLight, minimum width=1.5cm] (p)
            at (0,0) {\textbf{Prefix}\\{\scriptsize 0--1 B}};
        \node[bidang, fill=obeBlueLight, minimum width=2.6cm]
            (o) at (2.6,0) {\textbf{Opcode}\\{\scriptsize 1--2 B}};
        \node[bidang, fill=obeTealLight, minimum width=2.6cm]
            (m) at (5.8,0) {\textbf{ModR/M}\\{\scriptsize 0--1 B}};
        \node[bidang, fill=obeGreenLight, minimum width=2.6cm]
            (d) at (8.9,0) {\textbf{Displacement}\\{\scriptsize 0--2 B}};
        \node[bidang, fill=obeMaroonLight, minimum width=2.2cm]
            (i) at (11.8,0) {\textbf{Immediate}\\{\scriptsize 0--2 B}};

        \node[ket, below=0.1cm of p] {opsional};
        \node[ket, below=0.1cm of o] {wajib};
        \node[ket, below=0.1cm of m] {bila ada operand};
        \node[ket, below=0.1cm of d] {bila alamat\\bergeser};
        \node[ket, below=0.1cm of i] {bila operand\\berupa angka};

        % --- baris kedua: contoh nyata ---
        \node[bidang, fill=obeBlueLight, minimum width=2.0cm]
            (o2) at (1.2,-2.9) {\hex{C7}};
        \node[bidang, fill=obeTealLight, minimum width=2.0cm]
            (m2) at (3.2,-2.9) {\hex{80}};
        \node[bidang, fill=obeGreenLight, minimum width=2.6cm]
            (d2) at (5.8,-2.9) {\hex{34 12}};
        \node[bidang, fill=obeMaroonLight, minimum width=2.6cm]
            (i2) at (8.4,-2.9) {\hex{78 56}};

        \draw[-{Stealth[length=2.2mm]}, obeGray, line width=0.7pt]
            (o.south) -- (o2.north);
        \draw[-{Stealth[length=2.2mm]}, obeGray, line width=0.7pt]
            (m.south) -- (m2.north);
        \draw[-{Stealth[length=2.2mm]}, obeGray, line width=0.7pt]
            (d.south) -- (d2.north);
        \draw[-{Stealth[length=2.2mm]}, obeGray, line width=0.7pt]
            (i.south) -- (i2.north);

        \node[ket, left=0.1cm of o2, align=right]
            {\code{MOV WORD PTR}\\\code{[BX+SI+1234h],}\\\code{5678h}};

        \node[ket] at (6.2,-4.2)
            {\hex{C7 80 34 12 78 56} --- enam byte, semua field terpakai};
    \end{tikzpicture}
    \caption{Susunan field instruksi 8086 dan contoh penerapannya
    pada \code{MOV WORD PTR [BX+SI+1234h], 5678h}.}
    \label{fig:format-instruksi}
\end{figure}

\subsection{Ukuran Instruksi Tidak Seragam}

Karena field bersifat pilihan, panjang instruksi 8086 berkisar antara
1 dan 6 byte. Intel membatasi panjang maksimum instruksi pada 6 byte,
dan angka inilah yang menjadi alasan antrean pengambilan awal
(\eng{prefetch queue}) pada 8086 berukuran 6 byte
\cite{intel1979}. Tabel~\ref{tab:panjang-instruksi} memperlihatkan
beberapa contoh nyata.

\begin{table}[H]
    \centering
    \caption{Variasi panjang instruksi 8086 dan field yang
    dipakainya.}
    \label{tab:panjang-instruksi}
    \small
    \begin{tabular}{@{}p{4.6cm}p{3.1cm}cp{3.6cm}@{}}
        \toprule
        \textbf{Instruksi} & \textbf{Kode mesin} &
        \textbf{Byte} & \textbf{Field yang dipakai} \\
        \midrule
        \lstasm|NOP| & \hex{90} & 1 & Opcode \\
        \addlinespace
        \lstasm|PUSH AX| & \hex{50} & 1 & Opcode \\
        \addlinespace
        \lstasm|MOV AL, 5| & \hex{B0 05} & 2 & Opcode, immediate \\
        \addlinespace
        \lstasm|MOV AX, BX| & \hex{89 D8} & 2 & Opcode, ModR/M \\
        \addlinespace
        \lstasm|MOV AX, 1234h| & \hex{B8 34 12} & 3 &
        Opcode, immediate 16 bit \\
        \addlinespace
        \lstasm|ADD AX, [SI+4]| & \hex{03 44 04} & 3 &
        Opcode, ModR/M, displacement \\
        \addlinespace
        \lstasm|MOV WORD PTR| & & & \\
        \lstasm|[BX+SI+1234h],| & \hex{C7 80 34} & & Opcode, ModR/M, \\
        \lstasm|5678h| & \hex{12 78 56} & 6 &
        displacement, immediate \\
        \bottomrule
    \end{tabular}
\end{table}

\subsection{Membaca Kode Mesin dari Kiri ke Kanan}

Byte kode mesin dieksekusi menurut urutan alamat, dari byte pertama
sampai terakhir. Namun nilai 16 bit di dalamnya disimpan dengan
urutan \term{little-endian}, yaitu byte bernilai rendah lebih
dahulu. Akibatnya nilai \hex{1234} pada \lstasm|MOV AX, 1234h|
tertulis sebagai \hex{34} lalu \hex{12}. Kesalahan membaca urutan ini
termasuk kekeliruan yang paling sering terjadi pada analisis kode
mesin \cite{carleton2011}.
